\begin{abstract}

细粒度图像分类要求模型在同一大类内部区分视觉差异很小的子类,典型任务包括鸟类、车型和宠物品种识别。与普通图像分类相比,该任务往往面临类间差异小、类内变化大的问题,在少样本条件下尤其容易出现特征不足和类别混淆。Oxford-IIIT Pets 数据集中的猫狗品种识别就具有这一特点: 不同品种在毛色、脸型和体态上只存在局部差别,而拍摄姿态、光照和遮挡又会进一步干扰分类。

视觉语言预训练模型 CLIP 为小样本图像分类提供了新的实现路径。它可以利用图像与文本的对齐能力完成零样本或少样本识别,但分类效果仍然明显依赖提示词设计。已有的 CoOp 通过学习上下文向量替代人工提示词,改善了下游任务适配能力; CoCoOp 进一步引入图像条件化提示,提高了不同样本之间的适应性。不过,这两类方法都没有直接区分易分类样本和难分类样本,训练时仍采用统一的损失权重。

针对上述问题,本文提出一种基于 CoCoOp 的难样本加权扩展方法,并在项目中以 DynamicPromptTrainer 训练器形式实现。该方法以 CLIP 为基础,沿用可学习上下文和图像条件提示机制生成动态提示,并进一步引入难度权重计算器,根据模型对真实类别的预测置信度和错误情况生成可学习的样本权重,使训练过程对困难样本给予更高关注。与此同时,本文采用随 shot 数变化的自适应训练轮数策略,以缓解低样本条件下训练不足和高样本条件下过拟合的问题。

本文在 Oxford-IIIT Pets 数据集上完成了 1-shot、2-shot、4-shot、8-shot 和 16-shot 五组实验。结果表明,所提方法在 2-shot、4-shot、8-shot 和 16-shot 设置下均优于 CoCoOp,对应准确率分别达到 85.9\%、89.5\%、89.2\% 和 89.8\%; 其中 16-shot 设置下取得全实验最高准确率 89.8\%。实验结果说明,将图像条件提示与难度感知加权结合起来,能够更稳定地提升 CLIP 在细粒度少样本分类任务中的表现。

\keywords{细粒度分类,CLIP,提示词学习,少样本学习,猫狗品种识别}

\end{abstract}

\begin{abstractEN}

Fine-grained image classification aims to distinguish subcategories within the same coarse category, where visual differences are subtle and intra-class variation is often large. Pet breed recognition on Oxford-IIIT Pets is a typical example: breeds may differ only in local traits such as fur pattern, face shape, or body proportion, while pose, lighting, and occlusion still introduce strong appearance changes. These factors make the task difficult under few-shot supervision.

CLIP provides a practical foundation for this setting because it aligns image and text representations in a shared space. Yet its performance is still sensitive to prompt design. CoOp improves task adaptation by replacing handcrafted context words with learnable vectors, and CoCoOp further introduces image-conditional prompts. However, these methods do not explicitly distinguish easy samples from hard ones during training, and all samples still contribute to the loss in a similar manner.

This thesis presents a hard-sample weighting extension built on top of CoCoOp for fine-grained cat and dog classification, implemented in the project through the trainer named DynamicPromptTrainer. Built on CLIP, the method keeps the learnable context and image-conditional prompting scheme, and further introduces a difficulty weight calculator to assign learnable sample weights according to prediction confidence and misclassification status. In addition, an adaptive epoch schedule is adopted so that lower-shot settings receive longer training while higher-shot settings avoid unnecessary overfitting.

Experiments are conducted on the Oxford-IIIT Pets dataset with five few-shot settings: 1-shot, 2-shot, 4-shot, 8-shot, and 16-shot. The proposed method achieves 85.9\%, 89.5\%, 89.2\%, and 89.8\% accuracy in the 2-shot, 4-shot, 8-shot, and 16-shot settings, respectively, outperforming CoCoOp in all four cases. The best result reaches 89.8\% under the 16-shot setting. These results show that combining image-conditioned prompts with difficulty-aware weighting is effective for improving CLIP-based fine-grained few-shot classification.

\keywordsEN{Fine-grained Classification, CLIP, Prompt Learning, Few-shot Learning, Pet Breed Recognition}
\end{abstractEN}
